\documentclass[11pt]{article}
\usepackage{manual_docs_style}

\begin{document}

\docTitle{Website}
\label{sec:website}

The project website is the public-facing entry point for TSENV. 
\begin{itemize}
    \item Make the benchmark understandable before a reader opens the paper or repository: what TSENV measures, why time-series exploration is difficult for agents, how the tasks are generated, what results were observed, and how to reproduce the evaluation.
    \item Enable people to download the results and trajectories
    \item Allow external evaluators to submit their evaluations
\end{itemize}
We draw inspiration from \code{TerminalBench} and \code{TeamBench} websites. 
The website is hosted at tsenv.github.io.
Benchmark data, public examples, and downloadable artifacts are hosted on the Hugging Face Hub.
The website should be both desktop- and mobile-friendly.

\subsection*{Where is stuff hosted and located}
GitHub Pages = public website, documentation, and static result tables.
GitHub Actions = regenerate JSON and deploy website.
Hugging Face Dataset = public benchmark data, samples, trajectories, and downloads.
GitHub PRs = official submission mechanism.
Where HF Dataset repo:
\begin{verbatim}
tsenv/tsenv-benchmark
  |-- summary.json
  |-- results.parquet
  |-- all_samples.parquet
  |-- examples/
  |-- trajectories/
  `-- README.md / dataset card
\end{verbatim}

\subsection*{Update mechanism}
\begin{verbatim}
new results committed / PR merged
  -> GitHub Action runs:
       python scripts/generate_website_data.py
       python scripts/validate_website_data.py
       npm run build
  -> deploy to GitHub Pages
\end{verbatim}




\section*{Website textual mockup}
\subsection*{Global style}

The intended style is minimalistic and simple. 
It should follow the design principle of `Perfection is not when there is nothing more to add, but there is nothing left to take away'.
White background, simple navigation, strong typography, rectangular figures, thin divider lines, and a blue link style consistent with ETH SIPLAB pages should be the standard.

The website goal is to explain what TSENV measures, how tasks are generated, what results were observed, how users can inspect examples, and how external evaluations can be submitted.

The website should use a restrained research-lab style:

\begin{itemize}
    \item \textbf{Layout.} Use a wide centered content column with generous margins and strong vertical spacing.
    \item \textbf{Typography.} Use large bold section titles, short descriptive subtitles, and compact body text.
    \item \textbf{Navigation.} Use a simple top header menu with text links: \texttt{Home}, \texttt{Results}, \texttt{Environments}, \texttt{Get Started}, and \texttt{Authors}. The website must be implemented as a single-page application with client-side routing. These links should update the browser URL and render the corresponding route without loading separate HTML documents or merely scrolling to sections within one long page.
    \item \textbf{Visuals.} Prefer rectangular plots and tables.
    \item \textbf{Color.} Use mostly black, white, and grey. Use blue for links and active menu items. Use orange only as a secondary accent.
    \item \textbf{Components.} Prefer thin rules, tables, and flat panels over heavy cards or shadows.
    \item \textbf{Section labels.} Do not use eyebrow labels or small uppercase pre-headings above page titles or section titles. 
    Each page should begin directly with its main title and short subtitle/description.
    Section hierarchy should be communicated through headings, spacing, and thin divider lines rather than additional decorative labels.
\end{itemize}

\subsection*{Website structure and routing}

The website must be implemented as a single-page application (SPA), not as a collection of separate static HTML pages and not as a single long scrolling page with anchor sections.
Each top-level navigation item should map to a client-side route handled by the SPA router:

\begin{itemize}
    \item \texttt{/}: Overview
    \item \texttt{/results/}: Results
    \item \texttt{/environments/}: Environments
    \item \texttt{/get-started/}: Get Started and Submit
    \item \texttt{/authors/}: Authors, Citation, and Contact
\end{itemize}

The website should not include a separate Explorer page or \texttt{/explorer/} route.
Generated examples or demo visuals, if needed, should appear on the Overview page only.

The SPA should also support result-detail routes for individual leaderboard entries:

\begin{itemize}
    \item \texttt{/results/}: main leaderboard table
    \item \texttt{/results/\textless submission-id\textgreater/}: detail page for one submitted evaluation
\end{itemize}

The navigation menu should link to these routes directly.
Route changes should preserve the SPA shell and should not trigger full-page reloads.
Internal anchor links may be used within a route for subsections, but the primary website structure must remain client-side routed.

\subsection*{Page 1: Overview}

The homepage should immediately explain TSENV, show the key benchmark statistics, and provide a compact example task visual before ending with the design principles.

\subsubsection*{Header}

The top navigation should contain the TSENV wordmark on the left and a horizontal menu:

\begin{center}
\texttt{Home \quad Results \quad Environments \quad Get Started \quad Authors \quad Paper (Arkiv) \quad Code}
\end{center}

The active page should be highlighted in blue.
The header should have a thin horizontal rule below it, not a large navigation box.

\subsubsection*{Hero}

The hero should be a single full-width introductory block within the main centered content column.
It should not use a two-column layout and should not place secondary panels, tables, statistics, benchmark-scope content, task-definition content, or design-principle content beside the title.
The hero should contain only the main title and the two-paragraph subtitle/description.
All other homepage content should appear below the hero.

The hero should contain the title:

\begin{quote}
\textbf{TSENV: Controllable Time-Series Exploration Benchmark}
\end{quote}

The subtitle should explain the benchmark in two short paragraphs:

\begin{quote}
A benchmark for evaluating whether tool-using agents can perform evidence-grounded analysis of multivariate time series.

Tasks are generated from physical simulators with known interventions.
Agents must combine textual descriptions, noisy observations, and optional labeled examples to identify what changed, or that no intervention occurred.
\end{quote}


\subsubsection*{Homepage content order}

Below the hero, the landing page should present the homepage sections in the following order:

\begin{enumerate}
  \item Statistic row
  \item Example task visual
  \item Design principles
\end{enumerate}

The \texttt{Design principles} section should appear near the bottom of the landing page, below the example task visual.

\subsubsection*{Example task visual}


The landing page should show a compact example task visual consisting of a full-width time-series plot with an intervention marker.
The example time-series plot should be rendered from a small static JSON file, \texttt{public/data/example\_plot.json}.
The JSON should include the intervention time, changed parameter, observed channels, candidate parameters, and decimated samples for the plotted channels.
The example task visual should use a single-column layout. 
The Task and Axis controls should appear above the plot.
The time-series plot should appear below the controls and should span the full available content width.
The prompt panel should appear below the plot, not beside it.
The prompt panel should contain the task prompt and optional context/examples text.

The example visual should include compact controls grouped into two labelled rows:

\begin{verbatim}
Task: [DIRECT] [CODE]
Axis: [NOISE] [CONTEXT] [EXAMPLES]
\end{verbatim}

\texttt{DIRECT} and \texttt{CODE} should behave as mutually exclusive task-mode controls.
Clicking \texttt{DIRECT} should show the direct-answer prompt, where the agent returns JSON predictions.
Clicking \texttt{CODE} should show the code-mode prompt, where the agent writes an executable \texttt{rule.py} file.

\texttt{NOISE}, \texttt{CONTEXT}, and \texttt{EXAMPLES} should behave as condition-axis toggles.
When \texttt{NOISE} is highlighted, the same example plot should be shown with visibly noisier observations while preserving the same intervention marker and underlying trajectory.
When \texttt{CONTEXT} is highlighted, the textual task description should include additional environment information, such as observed channels, candidate parameters, and relevant physical context.
When \texttt{CONTEXT} is not highlighted, the description should remain short and minimal.
When \texttt{EXAMPLES} is highlighted, the prompt should show that optional labeled examples are available to the agent; when \texttt{EXAMPLES} is not highlighted, no labeled examples should be shown.

The toggles should make the benchmark conditions understandable without adding a separate Controllable axes section to the homepage.
The interaction should remain simple: highlighted buttons indicate the currently active task mode and selected condition-axis variants.

An agent, identified by a robot icon, should pop up the message ``Gravity changed''.

After the example task visual, the homepage should include a section titled \texttt{Design principles}.

\paragraph{Key features}
\begin{itemize}
  \item \textbf{Physics-grounded generation.} Samples are produced by simulators rather than by static datasets, giving reliable intervention labels and controllable task variants.
  \item \textbf{Evidence-grounded agent evaluation.} Correct answers require inspecting the observed signals; textual context alone is not sufficient.
  \item \textbf{Controlled difficulty axes.} TSENV varies textual context, observation noise, labeled examples, and answer interface in matched experimental conditions.
\end{itemize}


\subsubsection*{Statistic row}

The homepage should include three compact statistics in the following exact order:

\begin{center}
\begin{tabular}{ll}
\hline
\textbf{Statistic} & \textbf{Meaning} \\
\hline
\textbf{3 physical simulators} & BallDrop, BounceBall, MassSlide \\
\textbf{2 task modes} & Direct and Code \\
\textbf{3 controllable axes} & noise, context, examples \\
\hline
\end{tabular}
\end{center}

The order should emphasize the benchmark structure first: simulators, then task modes, then controllable axes.


\subsection*{Page 2: Results}

The results page should follow a single leaderboard-table layout, similar in spirit to the Terminal-Bench leaderboard.
It should not show separate grouped result columns for \texttt{Direct} and \texttt{Code} on the main page.
Instead, the main results page should contain one sortable and filterable leaderboard table where each row corresponds to one submitted result.

\subsubsection*{Main leaderboard table}

The main results page should show a single leaderboard table with one row per submitted result.
The table should be the primary object on the page.

Recommended columns:

\begin{center}
\begin{tabular}{rllllll}
\hline
\textbf{Rank} & \textbf{Agent} & \textbf{Model} & \textbf{Task} & \textbf{Date} & \textbf{Submitter} & \textbf{Score} \\
\hline
1 & \texttt{gpt-5.5-agent} & \texttt{gpt-5.5} & Direct & 2026-05-01 & TSENV & 0.933 \\
2 & \texttt{claude-agent} & \texttt{claude-opus-4.6} & Direct & 2026-05-01 & TSENV & 0.853 \\
3 & \texttt{gemini-agent} & \texttt{gemini-3.1-pro} & Code & 2026-05-01 & TSENV & 0.733 \\
\hline
\end{tabular}
\end{center}

The main leaderboard should use simple table styling with thin rules and no heavy dashboard cards.
The table should support sorting by rank, score, model, task, date, and submitter.

The \texttt{Task} column should distinguish \texttt{Direct} and \texttt{Code}.
Simulator, noise level, textual context, examples, and task should be available as filters.
Seed, split, and submitter may remain available as metadata, but should not appear as main filters.

\subsubsection*{Condition filters}

The page should include compact filters above the leaderboard table for:

\begin{itemize}
    \item simulator,
    \item noise level,
    \item textual context,
    \item examples,
    \item task.
\end{itemize}

The filters should be simple and readable, similar to a research benchmark leaderboard.
They should not make the page look like a heavy dashboard.


\subsubsection*{Clickable leaderboard rows}

Each leaderboard row should be clickable.
Clicking a row should navigate to a result-detail route rather than expanding a panel in place.

For example:

\begin{verbatim}
/results/<submission-id>/
\end{verbatim}

The result-detail page should show:

\begin{itemize}
    \item the selected agent, model, task, date, submitter, and score;
    \item the exact benchmark condition: simulator, noise, textual context, examples, seed, split, and task;
    \item cumulative tokens, Python calls, steps, and plots;
    \item per-environment results;
    \item per-seed results;
    \item links to trajectories and downloadable artifacts.
\end{itemize}

The browser URL should update when a user opens a result-detail page.
The route should remain inside the single-page application and should not trigger a full-page reload.


\subsection*{Page 3: Environments}
The page should have a grid layout where each environment is a square. 
When you click on the square you get to an environment-specific layout.
The top of the page should show a plot of the time series, followed by the description.
OBSERVED CHANNELS
height, velocity, contact impulse, time
CANDIDATE PARAMETERS
mass, drag coefficient, coefficient of restitution, gravity acceleration

For example with reference to BallDrop
\subsubsection*{BallDrop}
The BallDrop row should include:
\begin{itemize}
    \item a time-series plot with an intervention marker;
    \item a short description: vertical point-mass ball under gravity, quadratic drag, and restitution-based ground impacts;
    \item observed channels: height, velocity, contact impulse, and time;
    \item candidate parameters: mass, drag coefficient, restitution, and gravity.
\end{itemize}


\subsubsection*{Layout notes}

The environment page should be easy to scan vertically.
Plots should be large enough to be useful.
Observed channels and candidate parameters should be aligned in labeled rows, not buried inside prose.

\subsection*{Page 4: Get Started and Submit}

The get-started page should be practical.
It should contain a large copy-pasteable code block and a clear submission flow.

\subsubsection*{Run locally}

The quick-start section should contain:

\begin{verbatim}
# Clone and install
git clone <TSENV_REPOSITORY_URL>
cd tsENV
python -m venv env
source env/bin/activate
pip install -e .

# Launch local explorer
bash web_model_explorer/start.sh

# Generate question bundles
python workflows/generate/train_test_selection.py --model BallDrop
python workflows/generate/create_exam_questions_tsenv_cls_from_registry.py \
  --model BallDrop --row-slug <ROW_SLUG>

# Run one agentic evaluation
python workflows/rollout/question_run_orchestrator.py \
  --tasks-dir tsENV_questions \
  --model BallDrop \
  --agent-id <AGENT_ID> \
  --row-slug <ROW_SLUG>
\end{verbatim}

The code block should be visually prominent because this page is primarily a usage page.

\subsubsection*{Submission flow}

The submission flow should have four steps.
The third step should be labeled clearly as \textbf{Submit / Open PR}, based on the handwritten print annotation.

\begin{center}
\begin{tabular}{lp{0.68\linewidth}}
\hline
\textbf{Step} & \textbf{Description} \\
\hline
1. Run & Produce predictions, logs, and score artifacts \\
2. Validate & Check schema, date format, task counts, and condition metadata \\
3. Submit / Open PR & Upload results to a branch and open a pull request \\
4. Trace & Link each public score to generated artifacts \\
\hline
\end{tabular}
\end{center}

The page should also show the expected deliverables:

\begin{itemize}
    \item \texttt{results/*.json}
    \item \texttt{logs/}
    \item \texttt{artifacts/}
    \item \texttt{metadata.json}
    \item optional notes or README file
\end{itemize}

At the bottom, link back to the \texttt{Results} page for leaderboard structure and to the documentation for schema details.

\subsection*{Page 5: Authors, Citation, and Contact}

The final page should make the website feel complete, citable, and easy to contact.
It should include authors, citation information, resources, and contact links.
It should not include an FAQ section.

\subsubsection*{Authors}

The page should include a clear \textbf{Authors} section listing the project authors, affiliations, and optionally equal-contribution notes.

\subsubsection*{Citation}

The citation should be shown in a dark code block:

\begin{verbatim}
@inproceedings{tsenv2026,
  title={TSENV: Controllable Time-Series Exploration Benchmark for Agents},
  author={...},
  year={2026}
}
\end{verbatim}

If the final publication venue is not known yet, the citation block should either use a placeholder or say \texttt{Citation coming soon}.

\subsubsection*{Contact}

The final page should include a clear \textbf{Contact} section.
This was added based on the print annotation.

Recommended contact links:

\begin{center}
\texttt{Email \quad Repository \quad Affiliation}
\end{center}

The short description should be:

\begin{quote}
Questions, feedback, or collaboration.
\end{quote}

\subsubsection*{Resources}

The resources section should contain text links:

\begin{center}
\texttt{paper \quad code \quad dataset \quad results}
\end{center}

If one of these artifacts is unavailable, the link should be visible but disabled or marked as ``coming soon''.

\subsection*{Implementation checklist}

The website should pass the following checks before being considered complete:

\begin{itemize}
    \item No leftover \texttt{eda-bench}, \texttt{terminal-bench}, \texttt{example.com}, or placeholder repository text.
    \item The homepage explains TSENV within the first screen.
    \item The results page uses a single leaderboard table with one row per submitted result and a \texttt{Task} column distinguishing \texttt{Direct} and \texttt{Code}.
    \item The environments page lists observed channels and candidate parameters for all simulators.
    \item The submission page uses \textbf{Submit / Open PR} as the third step.
    \item The final page is titled Authors, Citation, and Contact, and does not include an FAQ section.
    \item The website is a single-page application with client-side routing; top-level navigation changes routes without loading separate HTML documents.
    \item The homepage hero is a single full-width block containing only the title and introductory description; no side panels, tables, statistics, benchmark-scope content, or task-definition content sit to the right of the hero title.
    \item The homepage statistic row uses this exact order: simulators, task modes, controllable axes.
    \item The example task visual groups controls into two rows: Task with \texttt{DIRECT} and \texttt{CODE}, and Axis with \texttt{NOISE}, \texttt{CONTEXT}, and \texttt{EXAMPLES}.
    \item All result values are generated from consolidated artifacts rather than manually copied.
\end{itemize}

\section{Putting the website online}

The website should be published in two phases:

\begin{enumerate}
    \item publish the current static single-page application first;
    \item replace mock data with generated benchmark artifacts.
\end{enumerate}

\subsection*{GitHub Pages host}

Use GitHub Pages as the public website host.
This matches the intended architecture: public website and static result tables on GitHub Pages, benchmark data and downloads on Hugging Face, submissions through GitHub pull requests, and deployment updates through GitHub Actions.

For the cleanest URL, create a GitHub repository named:

\begin{verbatim}
tsenv.github.io
\end{verbatim}

The static website files should live at the repository root:

\begin{verbatim}
tsenv.github.io/
  index.html
  404.html
  assets/
    app.js
    styles.css
  public/
    data/
      summary.json
      leaderboard.json
      environments.json
\end{verbatim}

The current static mockup already has the important static files:

\begin{verbatim}
index.html
404.html
assets/app.js
assets/styles.css
\end{verbatim}

\subsection*{Publishing the mockup}

From a local checkout of the website mockup:

\begin{verbatim}
unzip tsenv_website_mockup.zip
cd tsenv_mockup

git init
git branch -M main
git remote add origin git@github.com:tsenv/tsenv.github.io.git

git add .
git commit -m "Initial TSENV website"
git push -u origin main
\end{verbatim}

In GitHub, configure Pages to deploy through GitHub Actions:

\begin{verbatim}
Repository -> Settings -> Pages -> Build and deployment -> Source -> GitHub Actions
\end{verbatim}

\subsection*{Deployment workflow}

Create:

\begin{verbatim}
.github/workflows/deploy.yml
\end{verbatim}

For the current vanilla static mockup, use:

\begin{verbatim}
name: Deploy website

on:
  push:
    branches: [main]
  workflow_dispatch:

permissions:
  contents: read
  pages: write
  id-token: write

concurrency:
  group: pages
  cancel-in-progress: true

jobs:
  deploy:
    environment:
      name: github-pages
      url: ${{ steps.deployment.outputs.page_url }}

    runs-on: ubuntu-latest

    steps:
      - uses: actions/checkout@v4

      - uses: actions/configure-pages@v5

      - uses: actions/upload-pages-artifact@v3
        with:
          path: .

      - id: deployment
        uses: actions/deploy-pages@v4
\end{verbatim}

The \texttt{actions/deploy-pages} action deploys the Pages artifact uploaded by the workflow.

\subsection*{SPA deep links}

Because the website is a single-page application, routes such as these must load the SPA directly:

\begin{verbatim}
/results/
/results/gpt-55-agent-direct-20260501/
/environments/ball-drop/
\end{verbatim}

Keep \texttt{404.html} so direct links and browser reloads on nested routes still load the application.
The current mockup already includes a \texttt{404.html} redirect.
Later, improve this so nested-route reloads do not visibly show a ``404'' state before the SPA renders.

\subsection*{Generated website data}

The first deployed website may use mock data, but production result data should not stay hardcoded in \texttt{assets/app.js}.
Move leaderboard and environment data into generated JSON files:

\begin{verbatim}
public/data/leaderboard.json
public/data/summary.json
public/data/environments.json
public/data/example_plot.json
public/data/submissions/<submission-id>.json
\end{verbatim}

Then load the data from \texttt{assets/app.js} with \texttt{fetch()}:

\begin{verbatim}
const submissions = await fetch("/public/data/leaderboard.json")
  .then((response) => response.json());
\end{verbatim}

The website should not manually copy final result values.
All public result values should come from consolidated benchmark artifacts.
The homepage example plot should follow the same principle: \texttt{public/data/example\_plot.json} should be generated from a real BallDrop trajectory artifact and shipped as a small static JSON file.
The browser should fetch this JSON and render it directly; it should not load parquet files or depend on Hugging Face at runtime for the landing-page example.

\subsection*{Generation pipeline}

Add data generation scripts:

\begin{verbatim}
scripts/
  generate_website_data.py
  validate_website_data.py
\end{verbatim}

The deployment workflow should eventually generate and validate website data before publishing:

\begin{verbatim}
steps:
  - uses: actions/checkout@v4

  - uses: actions/setup-python@v5
    with:
      python-version: "3.11"

  - run: pip install -r requirements.txt

  - run: python scripts/generate_website_data.py

  - run: python scripts/validate_website_data.py

  - uses: actions/configure-pages@v5

  - uses: actions/upload-pages-artifact@v3
    with:
      path: .

  - id: deployment
    uses: actions/deploy-pages@v4
\end{verbatim}

If the website later migrates to Vite or React, add the Node build steps and deploy the generated \texttt{dist/} directory instead of the repository root:

\begin{verbatim}
  - uses: actions/setup-node@v4
    with:
      node-version: 22

  - run: npm ci
  - run: npm run build
\end{verbatim}

\subsection*{Hugging Face downloads}

The website repository should not store large trajectories or benchmark artifacts.
GitHub Pages should stay lightweight: HTML, CSS, JavaScript, small JSON files, and static figures.

Leaderboard rows and download links should point to the Hugging Face dataset:

\begin{verbatim}
https://huggingface.co/datasets/tsenv/tsenv-benchmark/resolve/main/trajectories/...
https://huggingface.co/datasets/tsenv/tsenv-benchmark/resolve/main/results.parquet
\end{verbatim}

\subsection*{Launch checklist}

Before making the website public, replace placeholder content:

\begin{verbatim}
TSENV Project Team
Email coming soon
paper coming soon
https://arxiv.org/
placeholder citation
mock model names, if not real
\end{verbatim}

Also verify:

\begin{verbatim}
No eda-bench / terminal-bench / example.com leftovers
Homepage explains TSENV immediately
SPA routing works
Results are generated, not manually copied
Get Started has Submit / Open PR
Authors page has no FAQ
Mobile layout works
\end{verbatim}

\end{document}
